%==============================================================================
%  02-ai-paradigm.tex — Section 1: The AI Paradigm
%==============================================================================
\strategyPart{05 \textbar\ The AI Paradigm}{From Digitalization to AI-Enabled Reinvention}

Organizations do not fail at AI because the technology is immature. They fail
because they apply it with a mental model built for an earlier era. This
section defines the three waves of transformation and establishes the
philosophy that the remainder of this document operationalizes.

\section{The three waves of transformation}

\begin{center}
  \begin{tabular}{@{}p{0.30\linewidth}@{\hspace{0.6cm}}p{0.30\linewidth}@{\hspace{0.6cm}}p{0.30\linewidth}@{}}
    \aiCard{%
      \aiCardTitle{1 · Digitalization}\\[0.4em]
      Converting analog records, processes and transactions into digital form.
      The baseline of the modern enterprise.
    } &
    \aiCard{%
      \aiCardTitle{2 · AI Adoption}\\[0.4em]
      Embedding point solutions — forecasting, chatbots, vision systems — into
      existing workflows without redesigning them.
    } &
    \aiCard{%
      \aiCardTitle{3 · AI-Enabled Reinvention}\\[0.4em]
      Redesigning business models, operating models and culture around machine
      intelligence. This is the frontier.
    } \\
  \end{tabular}
\end{center}

Most organizations are trapped between wave two and wave three: they have
adopted AI tools, but their processes, incentives and structures were designed
for a pre-AI world [McKinsey 2023]. The value does not come from the tool; it
comes from the reinvention the tool enables.

\section{The macroeconomic case}

The scale of the opportunity is measurable:
\begin{itemize}
  \item AI could add up to \textbf{\$13~trillion} to global economic output by
        2030 — roughly 16\% of cumulative global GDP [McKinsey Global
        Institute 2018].
  \item Generative AI alone could add \textbf{\$2.6--\$4.4~trillion} per year
        across 63 use cases, and could automate 60--70\% of the activities that
        currently consume employees' time [McKinsey 2023].
  \item Adoption is accelerating: 65\% of organizations now regularly use
        generative AI in at least one function, up from roughly one-third ten
        months earlier, and 72\% use AI in at least one business function
        [McKinsey State of AI 2024].
\end{itemize}

The strategic implication is blunt: the prize is real, the clock is running,
and the majority of incumbents are not converting adoption into value.

\section{Problem-first versus technology-first}

The organizing philosophy of this document is:

\begin{center}
  \textbf{\large ``We do not implement AI for the organization; we solve
  organizational problems using AI.''}
\end{center}

The technology-first fallacy proceeds from capability to problem: ``we have
a large language model — where can we use it?'' The problem-first discipline
proceeds from pain point to capability: ``which high-value organizational
problem is poorly served today, and which AI capability — predictive,
optimization, generative, or vision — can solve it?'' The research is
unequivocal about the consequences of getting this backwards: seven of ten
companies report minimal or no impact from AI so far, and among organizations
making significant AI investments, two in five report no bottom-line impact
[MIT SMR--BCG 2019]. The survivors and winners are distinguished not by the
volume of AI experiments but by the discipline with which they attach every
experiment to a business problem and a measurable outcome.
